\documentclass[11pt,a4paper]{article}
\usepackage[margin=23mm]{geometry}
\usepackage{fontspec,amsmath,amssymb,booktabs,tabularx,xcolor,hyperref}
\setmainfont{DejaVu Serif}
\setsansfont{DejaVu Sans}
\hypersetup{colorlinks=true,linkcolor=blue,citecolor=blue,urlcolor=blue}
\setlength{\parindent}{0pt}\setlength{\parskip}{6pt}
\setlength{\emergencystretch}{3em}
\title{Coherence Returns and Phase-Remainder Bounds in the Native Tri-Octagon Map}
\author{Hilmir Frímann Halldórsson}
\date{31 October 2026\\Research draft v0.1 — for GPT/author review}
\begin{document}\maketitle
% Generated from content.py; supplied LaTeX source was not compiled in this task.
% REVIEW.pdf was typeset from the same content with existing ReportLab/mathtext.
\section*{Abstract}
The native Tri-Octagon map evolves three complex channel states through a nonlinear prestage and a simultaneous harmonic-three phase update. We distinguish return of raw coherence from return of the full state. On a single regular block, equality of raw coherence is equivalent to a unit common-phase remainder; equivariance then propagates the block without assuming global regularity. For one fixed rational coefficient tuple, we prove absorption into a smaller amplitude domain, a sharper amplitude bound for every finite unit-multiplier return, and a phase-sector exclusion. Independent review extends the established 30-degree sector to the closed 60-degree sector: a regular finite return that meets this sector must be a fixed state. A separate three-update chord estimate bounds the principal remainder angle by approximately 18.5474 degrees and excludes the specific quarter-phase remainder mechanism on the stated domain. The retained 60-start investigation produced 42 numerical coherence-period-one outcomes in seven groups, 18 solver failures and no genuine period-three candidate. These finite numerical outcomes do not exclude other dispersed-phase returns. Existing state and observer evolution remain intact; a correspondence to the author's prescribed cone-pair choreography has not been established.

\subsection*{Contribution in ordinary language}
A repeated pattern can hide a common complex phase change. The paper identifies exactly when that happens, limits how large a three-update phase change can be at the studied settings, and shows that a finite return cannot sustain distinct stages inside a sufficiently narrow phase sector. The numerical search supplied examples of simpler returns, not a three-stage orbit. These are conclusions about the specified mathematical map, not measurements of energy leaving a geometric opening.

\subsection*{Authorship, status and evidence}
Human research author: Hilmir Frímann Halldórsson. GPT provided mathematical direction and internal review; Codex performed computational work in the retained investigations, the present source reconciliation and independent proof review, and manuscript preparation. This draft has not received author approval or external peer review. It is not publication approval. No new native-map evaluation was performed for this manuscript.

The quantitative results apply to the exact tuple in Section 2. Prior results at other tuples are separated. References link to a local bibliography; source IDs and relative retained-evidence paths are resolved in EVIDENCE\_INDEX.md. The manuscript and evidence can be reviewed without running scientific code.

\clearpage
\section*{1. Motivation and existing evolution}
The author described two cone pairs passing through a three-station sequence, with quarter-turn legs and an intended return comparison. The construction is a prescribed geometric question: which existing state evolution, if any, could induce that particular sequence? The original discussion sketch is not used as a figure or as evidence. The author’s golden-ratio wording was illustrative, not a requirement to construct or fit a controller. The sequence remains a motivation rather than a trajectory generated by the native recurrence. \cite{D01} \cite{D02}

Several kinds of evolution already exist in the project. The native map changes complex amplitudes and phases at integer updates. The separate SRG transfer acts on a six-component complex state. Decoded channel vectors change within fixed face planes. The staged and placed uppercase Z observations depend on their specified observer settings, and the EMA variant carries memory. Prescribed cone progress and display playback supply still other indices. A missing correspondence to cone-body pose does not imply that these earlier evolutions are absent. \cite{DYN} \cite{OBS} \cite{Z} \cite{BRIDGE}

\par\textit{Table 1. Distinct evolution and observation indices. Definitions and inherited context, not equivalences.}
\begingroup\small
\begin{tabularx}{\textwidth}{XXX}
\toprule
Object & Index or input & What is established\\\midrule
Native state z & Integer update n & Ordered nonlinear map F; amplitude and phase evolution\\
SRG state & Transfer count & Linear transfer U; separate extraction/initialization\\
Decoded vectors & Current native state & State-dependent vectors in fixed planes\\
Staged/placed Z & State plus observer sector/time & Specified readout; frozen and clocked cases differ\\
EMA Z & State and prior memory & Memory-dependent observation\\
Cone pairs & Prescribed leg progress and branch & Kinematic construction, not an F-orbit\\
Display & Playback time & Presentation setting, not a recurrence law\\
\bottomrule\end{tabularx}\endgroup\par
The focused question is therefore a return question. Can raw coherence have three distinct stages and repeat while the full state accumulates a common phase? If so, would a proposed readout turn that phase into the author’s geometric return? The first question can be formulated exactly without imposing the answer to the second. No fitted parameter, new time variable, new force law or state normalization is introduced. \cite{D02} \cite{D03}

Sections 3–6 give exact deductions in the NCM series; Section 5 contains the newly reviewed strengthening. Section 7 consolidates the completed numerical investigation. Sections 8–9 and Appendix B delimit inherited and kinematic comparisons. External novelty has not been assessed. The general return argument is elementary; the contribution is its precise application, quantitative bounds and comparison with retained evidence.

\clearpage
\section*{2. Native map, domains and coherence}
Let $z=(z_0,z_1,z_2)\in\mathbb{C}^3$. All coefficients are fixed real numbers. The graph matrix $L_3$ has diagonal entries $-2$ and off-diagonal entries $+1$; its spectrum is $0,-3,-3$. The ideal native update specified by the source is \cite{DYN}

\begin{equation}V_i(z)=z_i+\epsilon z_i(k_i-|z_i|^2)+g\sum_{j\neq i}(z_j-z_i).\end{equation}
\begin{equation}\delta_i(z)=\lambda\sum_{j\neq i}\sin\left(3(\operatorname{Arg}_0 V_j-\operatorname{Arg}_0 V_i)\right).\end{equation}
\begin{equation}F_i(z)=V_i(z)e^{\mathrm{i}\delta_i(z)},\qquad i=0,1,2.\end{equation}
Every kick uses the same prestage $V$. The convention is $\operatorname{Arg}_0(0)=0$; zero prestage outputs stay zero. The implementation bypasses phase extraction when $\lambda=0$. The exact ideal formulas, rather than floating-point implementation identities, are the objects of the proofs. The native source contains no added clock or amplitude normalization.

A regular $p$-block has $V_i(F^n z)\neq0$ for all three components and $0\leq n<p$. Regularity is a finite-block hypothesis. A zero prestage can break common-phase equivariance of the phase convention, because its assigned argument does not rotate with the state. We do not assume global equivariance at such zeros. \cite{OBS}

\begin{equation}H(z)=zz^\dagger,\qquad \operatorname{tr}H(z)=\|z\|_2^2,\qquad B_R=\{z:\max_i|z_i|\leq R\}.\end{equation}
Raw coherence $H$ retains total squared norm. It is not the normalized projector $H/\operatorname{tr}H$. The component-amplitude set $B_R$ is a polydisc in state space, unrelated to a spatial enclosing sphere. Common rephasing leaves $H$ unchanged, but that symmetry operation is not by itself one native update or a spatial rotation.

All quantitative claims in Sections 4–7 use exactly the following tuple, inherited from the retained study and coherent-relaxation source. Other parameter choices appear only in separated comparisons. \cite{COH} \cite{D03}

\begin{equation}\epsilon=\frac{1}{10},\quad g=\frac{1}{60},\quad\lambda=\frac{1}{300},\quad k=\left(\frac{999}{1000},1,\frac{1001}{1000}\right).\end{equation}
The completed investigation targeted $z\neq0$ in $B_3$, a regular three-block, and $F^3z=\mu z$ with $|\mu|=1$, where $H(z)$, $H(Fz)$ and $H(F^2z)$ are pairwise distinct. The phase remainder $\mu$ was unknown throughout. In particular, $\mu=\pm\mathrm{i}$ was not imposed on the solve.

\clearpage
\section*{3. A regular coherence return and its continuation}
\subsection*{Theorem 1. Return criterion and repeated blocks}
Fix arbitrary real coefficients in the native map. Let $z\neq0$ have a regular $p$-block. Then $H(F^p z)=H(z)$ if and only if there is a unique $\mu$ with $|\mu|=1$ and $F^p z=\mu z$. If this holds, all repeated blocks are regular and

\begin{equation}F^{mp+r}z=\mu^m F^r z,\qquad m\geq0,\quad 0\leq r<p.\end{equation}
This deduction in the NCM series uses the inherited regular phase-equivariance property. It is an elementary equivariant-map consequence, not a claim of a new universal mathematical principle. \cite{D02} \cite{OBS}

\textbf{Proof.} Put $w=F^p z$. Equality of the nonzero rank-one matrices $ww^\dagger$ and $zz^\dagger$ gives the same one-dimensional image, so $w=\mu z$. Taking traces gives $|\mu|=1$. Since $z\neq0$, the scalar is unique. The converse follows by substitution. The multiplier can be recovered without choosing a component:

\begin{equation}\mu=\frac{z^\dagger F^p z}{z^\dagger z}.\end{equation}
For a unit scalar $\alpha$, the prestage satisfies $V(\alpha x)=\alpha V(x)$ for every $x$, since amplitudes are unchanged. At a nonzero prestage, each principal argument changes by the common angle of $\alpha$ plus an integer multiple of $2\pi$. Differences in the harmonic-three sine therefore change by integer multiples of $6\pi$, leaving the kicks unchanged. Thus $F(\alpha x)=\alpha F(x)$ at a regular update. The prestage zero pattern also remains unchanged under multiplication by $\alpha$.

Apply this successively to the known $p$ updates. Starting at $F^p z=\mu z$, the next prestages are $\mu$ times the corresponding known prestages, hence nonzero; the next states are $\mu$ times the known states. This constructs another regular block. Induction constructs all later blocks and proves the formula. No assertion that an arbitrary regular state stays regular was used. This repeated orbit is regular because it is a rotated copy of its known block.

\subsection*{Corollary 1. Minimal full-state period}
Suppose $p$ is the minimal positive coherence period. If $\mu$ has exact finite order $q$, the full-state minimal period is $pq$. If $\mu$ is not a root of unity, there is no finite full-state return.

\textbf{Proof.} Write a proposed full return time as $N=ap+r$, with $0\leq r<p$. Theorem 1 gives $F^N z=\mu^aF^r z$. Equality with $z$ implies $H(F^r z)=H(z)$, so minimality forces $r=0$. The remaining condition is $\mu^a=1$. This proves both assertions. A numerical angle does not establish an exact finite order; rounding it to a rational fraction of a turn would change the problem.

\clearpage
\section*{3.1. Why normalized projectors are insufficient}
If only $H(w)/\operatorname{tr}H(w)=H(z)/\operatorname{tr}H(z)$ is known for nonzero states, the rank-one argument gives $w=az$ with an arbitrary nonzero complex scalar $a$. It does not give $|a|=1$. The native map is not homogeneous under amplitude changes: the cubic prestage changes its relative scale. Consequently the repeated-block conclusion cannot be obtained by replacing $\mu$ with such an $a$. \cite{D02}

The retained exact counterexample uses the separate aligned control tuple $k=(1,1,1)$, with $\epsilon=1/10$, $g=1/60$ and $\lambda=1/300$. It is not an example at the unequal-coefficient tuple of Section 2. For $z=r(1,1,1)$ and the positive amplitudes below, coupling vanishes and the phase kick is zero. The retained scalar expression is

\begin{equation}r^+=r\left[1+\frac{1-r^2}{10}\right].\end{equation}
\begin{equation}r_0=\frac{1}{2},\qquad r_1=\frac{43}{80},\qquad r_2=\frac{2947693}{5120000}.\end{equation}
All three normalized projectors agree. Yet repeating the first scale factor gives a different second amplitude:

\begin{equation}\left(\frac{r_1}{r_0}\right)r_1=\frac{1849}{3200}\neq\frac{2947693}{5120000}=r_2.\end{equation}
These are retained exact fractions, with their arithmetic rechecked symbolically for the manuscript. No native update was evaluated to produce a new example. The counterexample isolates the missing norm information: raw coherence has the strength required by Theorem 1, whereas normalization discards it.

\subsection*{Regularity and the reach of the criterion}
The algebraic rank-one equivalence itself does not need regularity. Regularity is needed when propagating the phase remainder through the actual ordered map and concluding that later blocks repeat. At a zero prestage, assigning zero to its phase is part of the model, not a removable numerical inconvenience. A theorem about a smoother replacement would not be a theorem about this map.

A raw coherence period also does not identify a pose period. A readout depending only on $H$ repeats with the coherence block, while a phase-sensitive readout can behave differently. To infer body motion, one would need a specified correspondence from the native state or its readouts to body coordinates, together with a demonstration that the correspondence respects the native update. Neither a common phase symmetry nor a projector return supplies that correspondence. Section 8 uses the existing cubic innovation as an example of phase sensitivity without assigning it a new mechanical meaning.

\clearpage
\section*{4. Amplitude absorption and returning-block bounds}
\subsection*{Lemma 2. Absorption into the smaller domain}
At the exact tuple of Section 2, $F(B_3)\subset\operatorname{int}B_{3/2}$. Both $B_{3/2}$ and $B_{\sqrt{K}}$ are forward invariant, where $K=1001/1000$. These conclusions do not require regularity. \cite{D03}

\textbf{Proof.} Set $a_i=1+\epsilon k_i-2g$ and $A=\max_i a_i$. Then

\begin{equation}\frac{31997}{30000}\leq a_i\leq A=\frac{32003}{30000}.\end{equation}
\begin{equation}a_i-\frac{r^2}{10}\geq\frac{4997}{30000}>0\qquad(0\leq r\leq3).\end{equation}
The phase update preserves component magnitudes, including at zero. For $z\in B_3$ and $r=|z_i|$, positivity of the self coefficient gives

\begin{equation}|F_i(z)|=|V_i(z)|\leq Ar-\frac{r^3}{10}+\frac{1}{10}.\end{equation}
The polynomial $h(r)=Ar-r^3/10$ reaches its nonnegative maximum at $r=\sqrt{10A/3}$. Its peak squared is $40A^3/27<49/25$. Thus $h(r)<7/5$ and every output modulus is strictly below $3/2$. This proves absorption without cancellation assumptions or differentiating an argument at zero.

For $R\leq3/2$, the upper envelope $h$ is increasing on $[0,R]$ because $h'(r)\geq11753/30000>0$. For $z\in B_R$, the sharper neighbor estimate gives

\begin{equation}|F_i(z)|\leq AR-\frac{R^3}{10}+2gR=T(R),\quad T(R)=R\left[1+\frac{K-R^2}{10}\right].\end{equation}
Now $T(3/2)<3/2$ and $T(\sqrt{K})=\sqrt{K}$, proving both invariances. The minimum derivative of an individual polynomial $a_i r-r^3/10$ on $[0,3/2]$ is instead $11747/30000$. These two positive derivative bounds are different and are not interchanged.

\subsection*{Proposition 3. Bound for any finite unit-multiplier return}
If $z\neq0$ belongs to $B_3$ and $F^p z=\mu z$ for $p\geq1$, $|\mu|=1$, every state in this finite block has component modulus at most $\sqrt{K}$. No regularity assumption is needed.

\textbf{Proof.} Absorption and invariance place updates 1 through $p$ in $\operatorname{int}B_{3/2}$. Since the last state has the same component moduli as $z$, the initial state does too. Let $M>0$ be the largest component modulus over stages 0 through $p-1$. Outputs over stages 1 through $p$ have the same maximum: the last stage replaces the initial one with identical moduli. The increasing envelope gives $M\leq T(M)$. If $M>\sqrt{K}$, then $T(M)<M$, a contradiction. This finite-block maximum argument uses neither global phase equivariance nor an infinite regular orbit. It supplies no positive lower amplitude bound and no classification outside $B_3$.

\clearpage
\section*{5. Closed 60-degree phase-sector exclusion}
\subsection*{Theorem 4. A regular finite return in the sector is fixed}
At the tuple of Section 2, let $z\neq0$ lie in $B_3$ and have a regular $p$-block satisfying $F^p z=\mu z$, $|\mu|=1$. If any stage of that block has its three component phases in a closed ordinary arc of width at most $\pi/3$, then $z$ is a fixed state and $\mu=1$.

This is the newly reviewed strengthening of the established 30-degree result. The disposition is ACCEPTED. The independent proof review is supplied in proof\_review/SECTOR\_60\_REVIEW.md; the original evidence is unchanged. Ordinary phases are modulo a full turn, without identifying phases separated by one third of a turn. This complex-phase-sector width is unrelated to the 60-degree spatial cone-axis midpoint. \cite{D03} \cite{REV03}

\subsection*{Proof, part I: positive prestage mixing}
Proposition 3 places the entire returning block in $B_{\sqrt{K}}$. Regularity makes every output component nonzero, and $F^p z=\mu z$ makes the initial components nonzero too. Theorem 1 permits a cyclic start at the sector stage without changing the multiplier.

\begin{equation}V_i=b_i z_i+g\sum_{j\neq i}z_j,\quad b_i=1+\epsilon k_i-2g-\epsilon|z_i|^2\geq\frac{14497}{15000}>0.\end{equation}
Choose consistent lifts $\theta_i$ in their shortest interval $[a,b]$, with $D=b-a\leq\pi/3$. If $D>0$, every prestage is a positive combination of all three nonzero rays. After rotation by $-a$, each sum has positive imaginary part, because some input is strictly above $a$ and every angle difference is less than $\pi$. After rotation by $-b$, it has negative imaginary part. Positivity also confines it to the pointed cone between the boundary rays, so it cannot be zero. Every prestage argument lies strictly in $(a,b)$; their enclosing width is strictly smaller than $D$. This is strict shrinkage for the particular state, not a uniform contraction rate.

\subsection*{Proof, part II: the actual simultaneous phase update}
Choose prestage lifts $\phi_i$ in that interval. For $d_{ij}=\phi_j-\phi_i$, put $a_{ij}=\lambda\sin(3d_{ij})/d_{ij}$ when $d_{ij}\neq0$, and $a_{ij}=3\lambda$ otherwise. On $|d_{ij}|\leq\pi/3$, the sine ratio is nonnegative, since $3d_{ij}\in[-\pi,\pi]$. The bound $|\sin x|\leq|x|$ supplies its upper bound:

\begin{equation}0\leq a_{ij}\leq3\lambda=\frac{1}{100},\qquad1-\sum_{j\neq i}a_{ij}\geq\frac{49}{50}.\end{equation}
\begin{equation}\phi_i^+=\left(1-\sum_{j\neq i}a_{ij}\right)\phi_i+\sum_{j\neq i}a_{ij}\phi_j.\end{equation}
Expanding this convex combination gives exactly the simultaneous kick. At $d=\pm\pi/3$ its weight is zero; at $d=0$ its continuous value is $1/100$. Thus the synchronizer is nonexpanding and need not itself strictly shrink. Changing a lift by an integer multiple of $2\pi$ changes each harmonic sine argument by an integer multiple of $6\pi$. The kicks and complex outputs are unchanged, so principal-argument wraps do not invalidate this computation.

\clearpage
\section*{5.1. Return contradiction and aligned descent}
\subsection*{Proof, part III: a positive-width return is impossible}
The full update strictly reduces positive sector width: the prestage strictly narrows it and the synchronizer stays in that narrower interval. Later widths remain at most $\pi/3$. For rays contained in an arc shorter than $\pi$, their shortest enclosing interval has their shortest circular arc width. That quantity is unchanged by a common unit multiplier. A finite return up to $\mu$ must therefore restore the initial width. Strict decrease contradicts this when the initial width is positive, including if it becomes zero during the block.

This uses only the known finite block and its rotated repetitions from Theorem 1. It does not assert that all regular states form an invariant domain. It requires no numerical lower cutoff on amplitudes and no normalization. The remaining possibility is an aligned state of width zero.

\subsection*{Proof, part IV: a finite aligned return is fixed}
Write the aligned state as $z=e^{\mathrm{i}\theta}r$ with all $r_i>0$. Positive prestage coefficients preserve its common phase, and all kicks vanish. Hence a finite aligned unit-multiplier return has $\mu=1$. The real amplitude update is $r^+=r-\nabla W(r)$ for the inherited potential \cite{D1}

\begin{equation}W(r)=\epsilon\sum_i\left(\frac{r_i^4}{4}-\frac{k_i r_i^2}{2}\right)+\frac{g}{2}\sum_{i<j}(r_i-r_j)^2.\end{equation}
The cube $[0,\sqrt{K}]^3$ is invariant. Positivity gives nonnegative outputs; the increasing envelope bounds them by $T(\sqrt{K})=\sqrt{K}$. The cube is convex, so the full segment joining successive amplitudes lies in it. The Hessian is the diagonal matrix with entries $\epsilon(3r_i^2-k_i)$ plus $-gL_3$. Its largest eigenvalue is bounded above by

\begin{equation}L=\epsilon(3K-k_{\min})+3g=\frac{313}{1250}<2.\end{equation}
Taylor’s upper estimate along that segment, with displacement $-\nabla W(r)$, gives

\begin{equation}W(r^+)\leq W(r)-\left(1-\frac{L}{2}\right)\|\nabla W(r)\|_2^2.\end{equation}
\begin{equation}1-\frac{L}{2}=\frac{2187}{2500}>0.\end{equation}
Summing over the finite aligned return forces every gradient to vanish, so every aligned stage is the same fixed state. If the sector was met at a cyclic stage, follow that fixed stage through the repeated block back to the original start. The common multiplier is one, so the original state is that fixed state too. This uses no inverse-map assumption and completes the proof.

The potential is used only for aligned amplitude dynamics. No descent claim is made for arbitrary dispersed-phase complex trajectories. The theorem excludes a nontrivial finite return meeting this sector; it is not a global convergence theorem. States whose phases remain more dispersed fall outside the argument. The established 30-degree theorem follows as a subset of the accepted sector.

\clearpage
\section*{6. A necessary three-update phase-remainder bound}
\subsection*{Proposition 5. Chord bound and quarter-phase exclusion}
At the fixed tuple, a nonzero regular three-block from $B_3$ satisfying $F^3z=\mu z$, $|\mu|=1$, obeys

\begin{equation}|\mu-1|\leq\frac{161151}{500000}=0.322302.\end{equation}
\begin{equation}|\operatorname{Arg}\mu|\leq2\arcsin\left(\frac{161151}{1000000}\right)\approx18.5474^\circ.\end{equation}
The angle is principal. In particular, $\mu=\pm\mathrm{i}$ is excluded. If coherence period is minimally three, an order-four remainder, and hence this particular full period-twelve mechanism, is impossible here. Other native twelve-cycles and other three-stage coherence returns are not excluded by this statement. \cite{D03}

\textbf{Proof.} Proposition 3 puts the block in $B_{\sqrt{K}}$. Write $V(z)-z=D_z z$ with $D_z=\epsilon\,\mathrm{diag}(k_i-|z_i|^2)+gL_3$. This matrix is Hermitian. Bounds on the diagonal terms and the spectrum of $L_3$ give

\begin{equation}\operatorname{spec}(D_z)\subset\left[-\frac{251}{5000},\frac{1001}{10000}\right].\end{equation}
\begin{equation}\|V(z)-z\|_2\leq\frac{1001}{10000}\|z\|_2,\quad\|V(z)\|_2\leq\frac{11001}{10000}\|z\|_2.\end{equation}
Every kick satisfies $|\delta_i|\leq2\lambda=1/150$. The chord estimate $|e^{\mathrm{i}t}-1|\leq|t|$ yields

\begin{equation}\|Fz-z\|_2\leq c\|z\|_2,\quad c=\frac{1001}{10000}+\frac{1}{150}\frac{11001}{10000}=\frac{53717}{500000}.\end{equation}
Choose the stage with maximum Euclidean norm as cyclic start. Theorem 1 shows this cyclic block is regular with the same multiplier: $F^3(F^s z)=\mu F^s z$. Common phase does not change norms. Denote the cyclic stages by $x_0,x_1,x_2,x_3=\mu x_0$. Telescoping gives

\begin{equation}|\mu-1|\|x_0\|_2\leq\sum_{n=0}^2\|x_{n+1}-x_n\|_2\leq3c\|x_0\|_2.\end{equation}
Divide by the nonzero norm. On the unit circle, $|\mu-1|=2|\sin(\operatorname{Arg}\mu/2)|$, giving the angle bound. Finally $|\pm\mathrm{i}-1|=\sqrt{2}>161151/500000$. The decimal angle merely formats the exact formula; it is not an interval certificate or fitted phase.

\clearpage
\section*{7. The frozen 60-start investigation}
\subsection*{Design and numerical screens}
This section reports retained DUAL-03 evidence only. The final phase\_b report, frozen plan, ledger, all 60 attempt records and regression summaries are supplied unchanged in the package. The preliminary missing-packet status was resolved before the search. No prior scientific script or native update was executed for this manuscript. \cite{D03}

The study used 54 prescribed polar starts and six prescribed Cartesian starts, including retained conjugate seeds, at the exact tuple of Section 2. The largest seed component was the gauge pivot, with lowest index on an exact tie. One common rephasing made it positive real without changing amplitude. This pivot stayed fixed throughout each solve; no chart restarts were used.

At 80 decimal digits, the eight real unknowns were the real and imaginary parts of $z$ and $\mu$. All six real return equations, the unit-circle equation and the imaginary-pivot gauge equation were retained. The initial phase estimate came from the three-step correlation without favoring a quarter phase. Newton directions used the analytic real Jacobian and LU. Return and gauge residual rows had a fixed per-start scale $\max(1,\|z_{\mathrm{initial}}\|_2)$; this did not normalize the state inside the map.

The line search tried 21 lengths $2^{-j}$, $0\leq j\leq20$, for the prescribed Armijo decrease. Trials outside $B_3$, with nonpositive real pivot, or nonfinite coordinates were rejected before a native call. Exact-zero prestages were refused without repair. Limits were 80 iteration passes per start and 120,000 discovery call attempts. Failed starts were not replaced.

\par\textit{Table 2. Frozen numerical screens. These are acceptance criteria, not existence certificates.}
\begingroup\small
\begin{tabularx}{\textwidth}{XX}
\toprule
Quantity & Retained criterion\\\midrule
Absolute return norm & $\|F^3z-\mu z\|_2\leq10^{-45}$\\
Relative return norm & $\|F^3z-\mu z\|_2/\|z\|_2\leq10^{-40}$\\
Unit-circle and gauge errors & Each at most $10^{-45}$\\
Zero/near-zero stop & State norm below $10^{-25}$\\
Post-convergence regularity & Minimum prestage modulus above $10^{-20}$\\
Post-convergence gauge & Real pivot above $10^{-25}$\\
Post-convergence amplitude & $\max_i|z_i|\leq\sqrt{1.001}+10^{-35}$\\
Lower period one & All H separations and one-step proportionality residual below $10^{-25}$\\
Genuine three-stage candidate & Every pairwise H separation above $10^{-12}$\\
\bottomrule\end{tabularx}\endgroup\par
H separations are Frobenius differences divided by $\|z_0\|_2^2$. Borderline cases are not promoted to candidates. Group matching allows common phase and cyclic starting stage, with relative discrepancy below $10^{-25}$; it does not identify permutations or conjugates.

\clearpage
\section*{7.1. Complete retained outcomes}
\par\textit{Table 3. Final classification of all 60 prescribed starts, extracted from the retained ledger.}
\begingroup\small
\begin{tabularx}{\textwidth}{XXX}
\toprule
Outcome & Count & Interpretation\\\midrule
Numerical coherence period one & 42 & Seven numerical groups; lower period\\
Line-search failure & 17 & No accepted return conclusion\\
Iteration-ceiling failure & 1 & S01; no accepted return conclusion\\
Genuine period-three candidate & 0 & None passed distinct-stage screen\\
Period-three certificate attempts & 0 & No candidate reached this stage\\
Validated period-three cycles & 0 & No such finding in the study\\
\bottomrule\end{tabularx}\endgroup\par
The seven numerical common-phase/cyclic-start groups are not certified unique roots, certified branches of the drift theorem, or basins of forward attraction. Root-finder convergence differs from attraction under the native map. Conjugate groups were kept distinct. \cite{D03}

\par\textit{Table 4. Retained representatives, rounded for display. Phase is the three-step remainder angle in radians; no exact order is inferred. Small values are not certified zero.}
\begingroup\small
\begin{tabularx}{\textwidth}{XXXXXX}
\toprule
Group & Seed & Starts & Phase (rad) & Min. prestage & Max. H sep.\\\midrule
G01 & S02 & 11 & 2.09793e-53 & 0.9996 & 2.05e-51\\
G02 & S04 & 10 & 3.36531e-07 & 0.704518 & 8.31e-78\\
G03 & S10 & 6 & -3.36531e-07 & 0.704518 & 5.4e-81\\
G04 & S13 & 7 & 0 & 0.159143 & 9.92e-65\\
G05 & S16 & 6 & 0 & 0.159571 & 2.68e-64\\
G06 & S19 & 1 & -6.99375e-05 & 0.000927912 & 1.07e-61\\
G07 & S52 & 1 & 0 & 0.159091 & 0\\
\bottomrule\end{tabularx}\endgroup\par
S01 illustrates the need for a relative criterion. Its retained final norm was about $4.7027133\times10^{-25}$ and its absolute return residual about $8.306139\times10^{-26}$, but its relative residual was about $0.1766244$. It failed; it was neither an accepted return nor an exact zero root. S37 also failed, with minimum prestage modulus about $7.771686\times10^{-6}$ and relative residual about $0.00251262$. A small prestage alone did not establish singularity or a return.

Every original attempt and failure is retained. The seven groups do not enumerate every possible lower-period state. Eighteen failures do not prove that their starting regions contain no roots. The lack of a genuine three-stage candidate in this finite local search leaves the dispersed-phase return question unresolved.

\clearpage
\section*{7.2. Verification ownership and limitations}
\par\textit{Table 5. Native evaluations in the original study versus this manuscript task.}
\begingroup\small
\begin{tabularx}{\textwidth}{XXX}
\toprule
Owner and category & Attempts & Completed\\\midrule
Original DUAL-03 discovery & 8,211 & 8,211\\
Original independent tests/regression & 151 & 150\\
Original total & 8,362 & 8,361\\
This manuscript preparation & 0 & 0\\
\bottomrule\end{tabularx}\endgroup\par
The original difference between attempts and completions was one intentional test refusal. The retained direct-formula comparison reevaluated three-blocks of the 42 converged attempts at 80 digits, accounting for 126 independent calls. Its reported maximum discrepancy was about 6.69022 × 10⁻⁸¹. This is a retained numerical regression result, not a new execution or an existence certificate. \cite{D03}

The ledger records 733 iteration passes, 691 Newton directions and 674 accepted steps, alongside 449 outside-domain trial rejections, 273 pivot rejections and 1,270 insufficient-decrease rejections. These describe solver behavior, not counts of dynamical trajectories or independent theorems.

Original Codex check categories were 25 Stage A checks, 87 pre-search checks and 68 post-search checks. The prior GPT review separately reported 25 and 68 replays and 27 additional algebra checks. Ownership and overlap matter: these cannot be added into a new certification total. This task performed file/data reconciliation, limited exact-algebra checks, written proof review and document QA. It did not replay those suites. \cite{REV03}

\subsection*{What follows, and what remains open}
The sector theorem is a universal exclusion under its finite-return hypotheses. The chord estimate is a universal necessary condition for a regular three-return on the admitted domain. The numerical investigation has a different logical status: outcomes of 60 fixed local solves. Neither failure nor a list of simpler roots proves absence of other three-stage returns.

A regular three-return in $B_3$ whose phases stay outside the closed $\pi/3$ sector and whose multiplier meets the chord bound remains unresolved. Near-zero or near-singular behavior cannot be certified from point residuals. Existence would need a justified certificate; global exclusion would need an argument covering the remaining domain. Neither is claimed or attempted here.

No continuation, extra starts, longer traces or memory experiments were performed for this draft. Original records remain authoritative, including their failures. They can be read as data without importing or executing supplied scientific scripts.

\clearpage
\section*{8. Existing readouts and common-phase sensitivity}
Raw coherence is common-phase invariant. The inherited coherence closure and frozen staged uppercase $Z$ readout share this invariance on their stated regular domains. A frozen readout holds observer sector/time and relevant settings fixed; advancing a clock is a different operation. A memory-dependent observation is not a function of current $H$ alone. \cite{OBS} \cite{Z}

The existing EMA innovation includes the cubic quantity denoted by $\mathcal{K}$ here to distinguish it from $K=1001/1000$: \cite{Z}

\begin{equation}\mathcal{K}(z)=z_0\overline{z_1}z_2,\qquad j(z)=\operatorname{Im}\mathcal{K}(z).\end{equation}
\begin{equation}\mathcal{K}(\mu z)=\mu^2\overline{\mu}\,\mathcal{K}(z)=\mu\mathcal{K}(z)\quad(|\mu|=1).\end{equation}
A regular block return from Theorem 1 therefore gives the conditional sequence

\begin{equation}j(F^{mp}z)=\operatorname{Im}\left(\mu^m\mathcal{K}(z)\right).\end{equation}
This consequence of an existing readout explains how identical coherence can coexist with a phase-sensitive scalar innovation. It is not a new motion law. It does not establish autonomous closure of $(H,j)$, a full EMA-memory period, energy flux, spin or cone pose. Nor does it show that the full state realizes any particular allowed multiplier. \cite{D02} \cite{REV02}

A changing decoded vector is an established geometric observation of the state. It is not by itself evidence that an attached rigid cone turns with that vector. A phase-insensitive readout cannot reveal the hidden common phase. A phase-sensitive one may reveal it, but body coordinates and a compatible evolution correspondence must still be specified and justified.

\subsection*{No quarter-phase trajectory is supplied}
A hypothetical $\mu=\mathrm{i}$ sequence could be written using the covariance of $\mathcal{K}$. It would describe an algebraic assumption, not a measured trajectory. At the studied tuple and domain, Proposition 5 already excludes a three-update return with that remainder. We therefore give the covariance formula without presenting a hypothetical quarter-phase table as a candidate orbit.

The author’s question about which opening shows energy leaving a core is not answered by these readouts. This investigation contains no calibrated energy field or flux law. What it supports is a distinction among evolving native state, repeated coherence, phase-sensitive scalar observation and externally prescribed pose. This preserves existing motion content while keeping the geometric interpretation open.

\clearpage
\section*{9. Relation to inherited native results}
The exclusions concern one fixed tuple and admitted domain. They do not replace the earlier native turning, relative-equilibrium, coherent-relaxation or two-cycle results, which use different hypotheses.

\subsection*{Coherent relaxation and transient turning}
The inherited coherent-relaxation theorem uses the same coefficients as Section 2 on the smaller domain of component amplitudes between $0.9$ and $1.1$ and ordinary phase width at most $\pi/12$. It establishes invariance and convergence to a unique aligned positive profile up to common phase. The retained turning entrance is outside that theorem’s domain, so the theorem does not automatically capture every turning trajectory. Theorem 4 concerns finite returns meeting a wider sector, not convergence of every state in that wider sector. \cite{COH}

\subsection*{D1 common-phase drift}
The drift source studies $(\epsilon,g,\lambda)=h(1,1/6,1/30)$, $0\leq h\leq1/4$, with $k=(1-\eta,1,1+\eta)$ locally in $\eta$ near the twisted reference with components $2^{-1/2}e^{2\pi\mathrm{i}j/3}$. Its local branch has common-phase drift and stationary coherence under its stated continuation hypotheses. It has saddle character, not a demonstrated forward-attracting three-stage orbit or cone pose. At $h=0$, the map is identity; an extended generator rate is not a nonzero single-update advance. The seven groups in Section 7 are not thereby certified as D1 branches. \cite{D1}

\subsection*{M3 entrance capture and a two-cycle}
M3 uses $\epsilon=1/20$, $g=1/5$, $k=(1,1.2208964704604097,6.35310346037241)$ and the exact critical phase strength plus $1/200$, as specified in its source. This differs from Section 2. The normalized Fourier entrance has squared norm three. Its certificate establishes entrance into a trapping box at update 301, a sufficient certified index, not necessarily first entry. Subsequent convergence is to one constant common rotation of a full complex two-cycle, with the stated axial/chiral sign alternation. The finite entrance trajectory is not claimed to become exactly periodic at update 301. \cite{M3}

That limiting cycle has full-state and coherence period two. A current-state readout sampled with a fixed integer stride cannot turn an exact two-cycle into a genuine three-stage sequence: an odd stride preserves at most two stages and an even stride selects one. This obstruction does not cover a clocked, memory-dependent or different state correspondence.

These inherited results demonstrate native evolution already present. This note adds a regular-return criterion, fixed-parameter limitations and a bounded numerical account. It is a focused companion to the turning/convergence paper, not a replacement.

\clearpage
\section*{10. Conclusions and bounded open questions}
Four conclusions remain separate. First, the model already contains native state evolution and geometric or memory-dependent observations. Second, a regular raw-coherence return has an exact common-phase remainder, controlling repeated blocks. At the studied tuple, finite unit-multiplier returns are amplitude bounded; any regular return meeting the closed 60-degree ordinary phase sector is fixed. A regular three-return also has the small remainder-angle bound, excluding its quarter-phase specialization.

Third, the retained finite search found no genuine three-stage candidate. Simpler outcomes and failures do not prove absence in the dispersed-phase domain. Fourth, no native-to-cone evolution correspondence was established. The prescribed sequence remains a geometric motivation, not a computed native motion or a disproved physical theory.

The open questions are correspondingly specific: whether a regular three-stage return exists in the admitted dispersed-phase region; whether any numerical lower-period group can be rigorously identified and certified; and which explicitly defined readout, if any, can respect a desired pose evolution. These are open statements, not instructions for another experiment. This draft fixes neither a phase order nor an observer clock to obtain a preferred choreography.

\subsection*{Contribution and placement}
The general criterion is an elementary symmetry consequence specialized to the regular domain and zero convention. The series-specific deductions are the finite-block amplitude bound, sector exclusion, remainder bound and combination with a complete retained numerical account. The 60-degree extension is newly reviewed, with an independent proof. No broad priority claim is made against external literature, which was not searched for this task.

A focused mathematical research note with its evidence package is the proportionate placement. Author and GPT review of claim wording, attribution, source placement and unresolved interpretation remains appropriate before any public submission. Publication decisions lie outside this task. The current artifact is a complete local review draft.

\subsection*{Reproducibility statement}
All reported native computations were completed in DUAL-03. This drafting task used Command Prompt and the existing torment environment for bounded hashing, JSON extraction, exact symbolic arithmetic and document generation. It made zero native F evaluations, zero old solver/checker/certificate executions and zero observer or pose evolution experiments. No Git commands, model changes, installations or publication actions were performed. Original review archives are included unchanged.

\clearpage
\section*{Appendix A. Exact constants and proof dependencies}
\par\textit{Table A1. Exact arithmetic supporting the proofs. Symbolic checks support identities; written proofs supply the universal arguments.}
\begingroup\small
\begin{tabularx}{\textwidth}{XXX}
\toprule
Quantity & Exact value or bound & Role\\\midrule
$a_{\min},a_{\max}$ & $31997/30000,\ 32003/30000$ & Prestage coefficients\\
Self coefficient on B3 & $\geq4997/30000$ & Positive mixing\\
Envelope derivative on [0,3/2] & $\geq11753/30000$ & Monotone envelope\\
Individual derivative minimum & $11747/30000$ & Distinct bound\\
Returning-block self coefficient & $\geq14497/15000$ & Strict sector mixing\\
Phase neighbor weight & $0\leq a_{ij}\leq1/100$ & Closed 60-degree sector\\
Phase diagonal weight & $\geq49/50$ & Convex phase update\\
Aligned Hessian upper bound & $313/1250$ & Descent on cube\\
Descent coefficient & $2187/2500$ & Aligned fixedness\\
Perturbation spectrum & $[-251/5000,1001/10000]$ & Norm estimate\\
One-step norm coefficient & $53717/500000$ & Telescoping\\
Three-step chord bound & $161151/500000$ & Remainder exclusion\\
\bottomrule\end{tabularx}\endgroup\par
For absorption one may instead use the looser envelope coefficient $27/25$. Its squared peak is $5832/3125$, below $49/25$ by $293/3125$. This retained alternative gives the same strict $3/2$ bound and agrees with the sharper coefficient in Section 4. No decimal approximation carries a proof step.

The new exact-algebra utility checks rational constants, endpoint and equal-phase identities, phase-branch invariance of harmonic sine, the convex-combination identity and retained counterexample fractions. It contains no native map implementation and imports no project module. It is not an interval certificate, orbit computation or finite substitute for the continuous-sector proof.

Dependencies run from exact map definition to regular equivariance and rank-one algebra; absorption and finite-block maximum; positive mixing and convex phase update; aligned potential descent; and the separate Hermitian/chord estimate. Numerical search evidence is not needed to establish these deductions. Conversely, the deductions do not convert a point residual into a certified orbit.

The claim matrix and before/after hashes are separate review artifacts. Hashes establish byte identity, not correctness or independent authentication of original executions.

\clearpage
\section*{Appendix B. Compact SRG and cone comparisons}
The inherited SRG transfer is $U=B\otimes A$, with a unitary two-channel factor and weighted Fourier cycle $A^3=D I$. Here $D=e^{-0.269}(0.382)^2$ lies strictly between zero and one. The raw norm after $3m$ transfers is multiplied by $D^m$, excluding nonzero raw periodic states of this transfer. This is about $U$, not native $F$. Normalized visual repetition would be a different criterion. \cite{BRIDGE} \cite{D02}

The genuine helicity extraction $\chi^\dagger\otimes I$ intertwines $U$ with $\lambda_\chi A$ under its eigenvector hypothesis. That linear extraction and the separate native initializer do not establish a conjugacy to $F$. A shared threefold pattern or matching dimensions does not equate their updates.

\par\textit{Table B1. Retained prescribed cone return counts, in repeated legs. Branches are retained choices, not native periods.}
\begingroup\small
\begin{tabularx}{\textwidth}{XXX}
\toprule
What is compared & Branch s=+ & Branch s=−\\\midrule
Unoriented pair axes & 3 & 3\\
Persistent end labels & 6 & 3\\
Independently turned full frames & 6 & 12\\
Rigid-frame comparator & 6 & 3\\
\bottomrule\end{tabularx}\endgroup\par
The prescribed station order is green A, B, C and orange B, C, A, with retained chart handedness and two unresolved route branches. For the negative branch, the three-leg marked residual is +90 degrees about green X and −90 degrees about orange Y. Frame markers diagnose orientation; they are not asserted physical features supplied by the author. Equal-progress independent turning gives a 60-degree unoriented pair-axis midpoint and a directed angle of 120 degrees; the rigid comparator stays perpendicular. These are chosen kinematics. \cite{D01}

For matched face transport, let $E_i$ be a three-by-two orthonormal tangent frame and $T_{ij}=E_iE_j^T$. Then $E_j^TE_j=I_2$ gives $T_{ij}T_{jk}=T_{ik}$. A closed loop is $E_iE_i^T$, identity on tangent plane $W_i$ but only a projector in ambient space. Fixed per-label gauges conjugate that identity and cannot supply a nontrivial quarter-roll. This rules out that particular matched-transport identification, not every state-dependent or nonlinear correspondence. \cite{D02}

These comparisons explain why recurrence, transport and prescribed pose cannot be interchanged merely because their geometry is related. They preserve the branch choices and geometric motivation without promoting either to a native motion law.

\clearpage
\section*{Appendix C. Review materials and build provenance}
START\_HERE.md gives the read order. CLAIM\_SOURCE\_MATRIX.md maps load-bearing results and tables to exact locators and evidence types. EVIDENCE\_INDEX.md resolves source IDs and archive identities. SOURCE\_MANIFEST.json records the bounded before/after source audit. Full DUAL-02 and DUAL-03 archives are retained unchanged under packet/NCM\_PAPER\_01\_Codex\_Packet\_v0.1/evidence/.

The numerical authority is the DUAL-03 archive’s final phase\_b/REPORT.md and its frozen plan, ledger, summary and 60 JSON attempts. Preliminary Stage A files remain history. Prior GPT reviews are internal review records, not original Windows executions. Current drafting checks have separate outputs and replace neither category.

One content source supplies the LaTeX file and this PDF. No local LaTeX compiler was available. This review PDF was typeset directly under torment using the existing bundled ReportLab library and matplotlib mathtext; no package or font was installed. Expressions are rendered at high resolution from the same source. The LaTeX source is supplied for editorial use but was not compiled or validated by a LaTeX engine here. A later TeX build may paginate differently.

MANUSCRIPT\_QA.md records render inspection and build limitations. COMMANDS.txt and CLOSEOUT.txt distinguish prior scientific work from current document work. The exact-algebra results and independent sector review are in proof\_review/. The output manifest excludes itself; the external receipt records ZIP identity and member verification without circular hashing.

No private sketch, generated infographic, font file, kernel copy or installation is included. Inherited papers are referenced by title, version and locator rather than repackaged. The bounded audit checks inspected source bytes, not the whole repository. It is not a producer attestation or publication approval.

\subsection*{Reference conventions}
The bibliography is restricted to inspected local project sources and retained reports. A research or publication edition is not described as externally peer reviewed. No external literature priority assessment was performed. Source IDs resolve through the relative evidence index so that absolute Windows paths are not the only references.

\clearpage
\begin{thebibliography}{REV03}
\bibitem{DYN} Tri-Octagon project source (2026). \emph{Native deterministic nonlinear map}. kernel\_physics/dynamics.py, lines 1–112; current bytes matched retained manifest. \href{../EVIDENCE_INDEX.md\#dyn}{Source record DYN}.
\bibitem{OBS} Hilmir Frímann Halldórsson (2026). \emph{Equivariant spatial observables}. Native Response Geometry, observables\_closure.tex; selected definitions and regular closure, lines 156–236. \href{../EVIDENCE_INDEX.md\#obs}{Source record OBS}.
\bibitem{COH} Hilmir Frímann Halldórsson (2026). \emph{Transient Collective Turning and Coherent Relaxation in the Native Triad}. Repository research edition v0.1, 28 October; fixed coefficients and coherent-domain theorem. \href{../EVIDENCE_INDEX.md\#coh}{Source record COH}.
\bibitem{D1} Hilmir Frímann Halldórsson (2026). \emph{Cubic Common-Phase Drift on a Chiral Branch of the Tri-Octagon Map}. Research publication v0.1, 6 October; model, local drift and amplitude potential. \href{../EVIDENCE_INDEX.md\#d1}{Source record D1}.
\bibitem{Z} Tri-Octagon project source (2026). \emph{Staged and EMA Z observers}. kernel\_physics/z\_manifold.py; observer and cubic innovation definitions. \href{../EVIDENCE_INDEX.md\#z}{Source record Z}.
\bibitem{M3} Hilmir Frímann Halldórsson (2026). \emph{Geometry-Attached Axial Observables, Period-Doubling, and Certified Entrance Capture in the Tri-Octagon Map}. Paper G v0.1.1, 4 October; retained M3 result and its parameter/entrance hypotheses. \href{../EVIDENCE_INDEX.md\#m3}{Source record M3}.
\bibitem{BRIDGE} Tri-Octagon project, prepared for Hilmir Frímann Halldórsson (2026). \emph{SRG Final Mathematical Bridge: From lens geometry to recurrence and known equations}. Publication edition v0.1, 27 October; transfer and extraction statements. \href{../EVIDENCE_INDEX.md\#bridge}{Source record BRIDGE}.
\bibitem{D01} Hilmir Frímann Halldórsson; Codex investigation (2026). \emph{NCM-DUAL-01: Dual Cone-Pair Quarter-Turn Sequence}. Completed internal report, context/NCM\_DUAL\_01\_REPORT.md; prescribed kinematics. \href{../EVIDENCE_INDEX.md\#d01}{Source record D01}.
\bibitem{D02} Hilmir Frímann Halldórsson; Codex investigation (2026). \emph{NCM-DUAL-02: Existing-Equation Correspondence and Return Mechanisms}. Completed internal report v0.1 and unchanged evidence archive; context/NCM\_DUAL\_02\_REPORT.md. \href{../EVIDENCE_INDEX.md\#d02}{Source record D02}.
\bibitem{D03} Hilmir Frímann Halldórsson; Codex investigation (2026). \emph{NCM-DUAL-03: Native Three-Update Coherence Returns at Fixed Parameters}. Final phase\_b report v0.1, frozen plan and all 60 attempt records; context/NCM\_DUAL\_03\_REPORT.md. \href{../EVIDENCE_INDEX.md\#d03}{Source record D03}.
\bibitem{REV02} GPT, internal mathematical reviewer (2026). \emph{Internal review of NCM-DUAL-02}. reviews/NCM\_DUAL\_02\_GPT\_REVIEW.md; phase-sensitive readout discussion. \href{../EVIDENCE_INDEX.md\#rev02}{Source record REV02}.
\bibitem{REV03} GPT, internal mathematical reviewer (2026). \emph{Internal review of NCM-DUAL-03}. reviews/NCM\_DUAL\_03\_GPT\_REVIEW.md, Section 5; sector proposal independently reviewed here. \href{../EVIDENCE_INDEX.md\#rev03}{Source record REV03}.
\end{thebibliography}
\end{document}
